\chapter{Getting Started}
\label{ch:getting-started}

This chapter takes you from a fresh clone to an LED blinking under your own Ada,
in a few commands. The later chapters explain the runtime, the HAL and the
language techniques; this one just gets you running.

\section{What you need}\index{Alire}

\begin{description}[style=nextline]
\item[An ESP32-S3 board.] Any devkit with the built-in USB-Serial-JTAG port (most
  of them). One USB cable is enough -- it carries the console \emph{and} the
  debugger.
\item[Alire (\texttt{alr}).] The Ada package manager fetches the toolchains the
  build needs: the \texttt{gnat\_xtensa\_esp32\_elf} cross-compiler, a native
  \texttt{gnat\_native} (for the host tools), and \texttt{gprbuild}. Install Alire,
  then make sure its toolchains are on your \texttt{PATH} for the shell you build
  from.
\item[Serial-port permission.] On Linux, add yourself to the \texttt{dialout}
  group (and install the udev rule for the USB-Serial-JTAG, see the debugging
  chapter) so you can open \texttt{/dev/ttyACM0} without \texttt{sudo}.
\end{description}

\section{Get the code}

The runtime depends on two git submodules (the \texttt{bb-runtimes} fork that
carries the ESP32-S3 bareboard runtime, and \texttt{xtensa-dynconfig}), so clone
\emph{recursively}:

\begin{lstlisting}[style=sh]
git clone --recurse-submodules \
    https://github.com/rowsail/ada_esp32s3.git
cd ada_esp32s3
#  (if you already cloned without submodules:)
git submodule update --init --recursive
\end{lstlisting}

\section{The \texttt{./x} dispatcher}\index{x helper script}

The repository ships with over sixty ready-to-run examples, and one script at
its root -- \texttt{./x} -- drives all of them \emph{from inside the clone}. (For
your \emph{own} application, kept outside the repo, the counterpart command is
\texttt{esp32-ada}, \S\ref{sec:esp32-ada}; the two share the same verbs.) List the
examples, then build, flash and watch one:

\begin{lstlisting}[style=sh]
./x list                                   # the available examples + their profiles
./x run esp32s3_gpio0_blink -p /dev/ttyACM0   # build + flash + serial monitor
\end{lstlisting}

\texttt{./x run} builds the example, packages it (the Ada \texttt{esp\_elf2image}
host tool), writes it to the board (the Ada \texttt{esp\_flash} tool, over the
chip's ROM bootloader), resets, and streams the console. The first build also
generates the runtime for the chosen profile, so it takes a little longer; later
ones are incremental. The blink example toggles an LED and prints over the
console -- your first Ada running bare-metal on the S3.

The pieces are also separate commands when you want them:

\begin{lstlisting}[style=sh]
./x build   esp32s3_gpio0_blink            # cross-compile + link + package app.bin
./x flash   esp32s3_gpio0_blink -p /dev/ttyACM0
./x monitor -p /dev/ttyACM0                # just the serial console (115200)
./x clean   esp32s3_gpio0_blink
\end{lstlisting}

The examples are meant to be \emph{read}, not just run: each opens with a header
that says what it demonstrates, how to build and run it, what the console should
print, and what hardware (if any) it needs; magic numbers are named, and the
``why'' is in the code rather than only in this book. A short house style at
\texttt{examples/STYLE.md} records the bar, with \texttt{esp32s3\_gpio0\_blink}
and \texttt{esp32s3\_gdma\_copy} as the models to copy; \texttt{./x new} (and
\texttt{esp32-ada init}, \S\ref{sec:esp32-ada}) scaffold a new project already in
that shape. So once a chapter here points you at an example, the example itself
carries the rest of the explanation.

\section{Choosing a profile}\index{runtime profile}\index{light-tasking profile}\index{embedded profile}\index{full profile}

Most examples build under more than one runtime profile (Chapter~\ref{ch:profiles}).
Pass \texttt{--profile} to pick one; \texttt{auto} (the default) uses the
example's own:

\begin{lstlisting}[style=sh]
./x run esp32s3_heartbeat --profile embedded -p /dev/ttyACM0
\end{lstlisting}

\texttt{light-tasking} is the smallest, \texttt{embedded} adds exceptions, RAII
handles, a heap and the full HAL (the usual choice), and \texttt{full} adds the
complete Ada tasking model.

\section{Board configuration}\index{board.ads}

Each project owns a \texttt{board.ads} (flash and PSRAM sizes) that sizes the image
and the second-stage bootloader. Edit it directly, or:

\begin{lstlisting}[style=sh]
./x config esp32s3_gpio0_blink show
./x config esp32s3_gpio0_blink flash-size 4MB
\end{lstlisting}

\section{Debugging}

On-chip debugging needs the pinned OpenOCD + S3 GDB, fetched once; then
\texttt{./x debug} (or \texttt{F5} in VS Code) halts at the boot entry
\texttt{app\_main} (the shared boot glue; these images have no C \texttt{main}):

\begin{lstlisting}[style=sh]
./x get-debug-tools
./x debug esp32s3_gpio0_blink              # build, flash, OpenOCD + GDB, break at app_main
\end{lstlisting}

The Debugging chapter (Chapter~\ref{ch:debug}) covers GDB, decoding a Guru
Meditation, and the editor integration in full.

\section{Your own project, out of tree: \texttt{esp32-ada}}\index{esp32-ada launcher}\index{scaffolding}
\label{sec:esp32-ada}

\texttt{./x} is for the examples that ship \emph{inside} the clone. For your own
application you do not edit an example or copy the runtime into your tree: you
treat this repository as an \emph{SDK} and scaffold a standalone project anywhere
on disk. The launcher that does this -- the out-of-tree counterpart of
\texttt{./x}, with the same verbs -- is \texttt{esp32-ada}.

It becomes available when you \emph{source} the SDK's \texttt{export.sh} once per
shell. That one step puts \texttt{esp32-ada} on your \texttt{PATH}, exports
\texttt{ESP32S3\_ADA\_SDK} (so a project folder anywhere can find the SDK), and
adds the runtime and HAL projects to \texttt{GPR\_PROJECT\_PATH} -- so both
\texttt{gprbuild} \emph{and} the Ada Language Server resolve
\texttt{with "esp32s3\_rts.gpr"} with no path baked into your project:

\begin{lstlisting}[style=sh]
. /path/to/ada_esp32s3/export.sh    # once per shell (add to ~/.bashrc to persist)
mkdir ~/myblink && cd ~/myblink
esp32-ada init                      # scaffold app.gpr, board.ads, src/main.adb
#  ... edit src/main.adb ...
esp32-ada run -p /dev/ttyACM0       # build + flash + monitor
\end{lstlisting}

Every \texttt{esp32-ada} verb mirrors the matching \texttt{./x} one; the only
difference is the target. \texttt{./x build esp32s3\_gpio0\_blink} names an example
inside the clone, while \texttt{esp32-ada build} operates on the project folder you
are standing in:

\begin{longtable}{p{0.46\linewidth} p{0.48\linewidth}}
\toprule
\textbf{Command} & \textbf{What it does} \\ \midrule
\endhead
\texttt{esp32-ada init [\textit{dir}]} & Scaffold a project in \textit{dir} (or the current folder). \\
\texttt{esp32-ada build [-P \textit{prof}]} & Build to \texttt{app.bin}. \textit{prof} = \texttt{light-tasking} (default), \texttt{embedded}, or \texttt{full}. \\
\texttt{esp32-ada flash [-p \textit{port}]} & Build if needed, then flash over the USB ROM bootloader. \\
\texttt{esp32-ada run [-p \textit{port}] [-P \textit{prof}]} & Build, flash, and open the serial monitor. \\
\texttt{esp32-ada monitor [-p \textit{port}]} & Just the serial console (115200). \\
\texttt{esp32-ada clean} & Remove this project's build artifacts. \\
\texttt{esp32-ada config [show | flash-size \textit{S} | psram-size \textit{S}]} & Show or set this project's flash / PSRAM size (its \texttt{board.ads}). \\
\texttt{esp32-ada debug [-p \textit{port}]} & On-chip debug (OpenOCD + GDB), halting at \texttt{Main}. \\
\texttt{esp32-ada kill-openocd} & Kill every OpenOCD (release captured USB-JTAG ports). \\
\texttt{esp32-ada install-ide} & Install the VS Code extension (committed \texttt{.vsix}). \\
\texttt{esp32-ada install-vim} & Symlink the Vim/Neovim plugin. \\ \midrule
\bottomrule
\end{longtable}

\noindent\texttt{-p}/\texttt{--port} defaults to \texttt{\$ESPPORT} or
\texttt{/dev/ttyACM0}; \texttt{-C \textit{dir}} runs as if from \textit{dir}.
\texttt{init} writes a complete, self-contained project:

\begin{lstlisting}[style=sh]
myblink/
  app.gpr            # withs esp32s3_rts.gpr + esp32s3_hal.gpr (resolved by name)
  board.ads          # this project's flash / PSRAM sizes
  src/main.adb       # your code (procedure Main; boots both cores, then idles)
  build.sh flash.sh  # thin shims into the SDK's shared bare-boot
  .vscode/           # ALS project + build/flash/run tasks driving esp32-ada
\end{lstlisting}

Nothing from the SDK is copied in -- \texttt{app.gpr} references the runtime and
HAL \emph{by name} through \texttt{GPR\_PROJECT\_PATH}, and the build shims call the
SDK's shared bare-boot through \texttt{\$ESP32S3\_ADA\_SDK}. Add another SDK library
the same way: \texttt{with "<name>.gpr";} in \texttt{app.gpr} for anything under the
SDK's \texttt{libs/}. The default profile is \texttt{light-tasking}; pass
\texttt{-P embedded} (or \texttt{full}) to build under another
(Chapter~\ref{ch:profiles}).

\subsection{Taking an example out of tree to edit}\index{examples}

The in-repo examples are wired to their location -- an Alire path-pin in
\texttt{alire.toml}, a \texttt{\textit{name}.gpr}, and build shims relative to
\texttt{examples/common/bare/} -- so copying a whole example directory elsewhere
will not build. Instead, scaffold a fresh project and bring the \emph{sources}
across; the scaffold already supplies the \texttt{app.gpr}, board config and build
shims that make them build:

\begin{lstlisting}[style=sh]
. /path/to/ada_esp32s3/export.sh
esp32-ada init ~/myblink && cd ~/myblink
SDK=$ESP32S3_ADA_SDK
cp "$SDK"/examples/esp32s3_gpio0_blink/src/*.ad? src/   # the Ada sources (incl. main.adb)
esp32-ada build                                         # gpio0_blink uses the default light-tasking
esp32-ada run -p /dev/ttyACM0
\end{lstlisting}

A few things make the copy build cleanly:

\begin{itemize}
\item \textbf{Only \texttt{src/} is yours.} \texttt{esp32-ada init} already wrote
  \texttt{app.gpr} (\texttt{Source\_Dirs => "src"}, \texttt{Main => "main.adb"}),
  so the example's \texttt{src/main.adb} (a \texttt{procedure Main}) and its other
  units drop straight in. Overwrite the scaffold's \texttt{src/main.adb}.
\item \textbf{Match the profile.} Most examples assume \texttt{embedded} (the HAL's
  RAII handles need it); the scaffold defaults to \texttt{light-tasking}. The
  example's profile is shown by \texttt{./x list} (and its \texttt{README}); pass
  it with \texttt{-P}.
\item \textbf{Bring board config and C natives if customized.} If the example sets
  a non-default flash/PSRAM size, copy its \texttt{board.ads} (or run
  \texttt{esp32-ada config}); if it has C natives, copy its \texttt{glue.c}.
\item \textbf{Extra libraries.} If the example \texttt{with}s more of the SDK's
  \texttt{libs/} (beyond the HAL), add the matching \texttt{with "<name>.gpr";} to
  \texttt{app.gpr}.
\end{itemize}

\cnote{There is no fixed ``project location'' and no \texttt{idf.py}: the SDK is
found through \texttt{\$ESP32S3\_ADA\_SDK} (set by \texttt{export.sh}), so your
project folder can live anywhere and stays portable -- check it into its own git
repo, with no SDK paths baked into any file.}

\section{Where to go next}

\begin{itemize}
\item The anatomy of an application and the runtime profiles --
  Chapters~\ref{ch:anatomy}--\ref{ch:build}.
\item The peripheral HAL -- start at Chapter~\ref{ch:hal}; each driver has worked
  examples you can run with \texttt{./x run}.
\item The Ada techniques behind it all -- the Ada Language Techniques and
  Talking to Hardware chapters (Chapters~\ref{ch:embedded}
  and~\ref{ch:hardware}), including a Terminology section (\S\ref{sec:terms})
  for unfamiliar terms.
\end{itemize}

To start your \emph{own} project rather than run an in-repo example, use the
\texttt{esp32-ada} launcher (\S\ref{sec:esp32-ada}), including the recipe for
lifting an example's sources out of the tree to build on as your own.
